\answer{ex:13:1}{(a)~$q=100^{1/99}\approx1.048$, pas $\approx4.8\%$; plage $\approx2.2\cdot10^3$ ($67$~dB). (b)~Pas de $22f_1$, soit $220\%$ à $10f_1$.}
\answer{ex:13:2}{$P_{\text{échec}}\approx1.8\cdot10^{-4}$ et $2.8\cdot10^{-16}$: environ $300$ échecs par jour pour $S=2$ et $5\cdot10^{-10}$ par jour pour $S=4$.}
\answer{ex:13:3}{(a)~$H(s)=eF_1T_c/(1+sT_c)^2$, filtre du second ordre de fréquence de coupure $1/(2\pi T_c)\approx3.2$~Hz. (b)~$r\approx22$~imp./s ($\approx20$ d'après le premier harmonique).}
\answer{ex:13:4}{(a)~Pour $Gd=1/e$, la racine $s=-1/d$ est double, et pour aucun autre $G$ la racine la plus à droite n'est plus à gauche. (b)~$d=30\ms$: $G\approx12$~s$^{-1}$, $\tau=30\ms$; $d=150\ms$: $G\approx2.5$~s$^{-1}$, $\tau=150\ms$; la constante de temps est égale au retard.}
\answer{ex:13:5}{(a)~$E\approx0.427$, $T\approx1.70\,\tau_a=0.68$~s (le modèle avec $\tau=20\ms$ donne $0.84$~s). (b)~L'équation ne contient pas $I$; le rythme existe pour $b>\beta-1$.}
\answer{ex:13:6}{Pour $b=0.8$, pas de rythme; pour $b=1.05$, $1.2$, $1.5$, $2$, $4$, $T\approx2.73$, $1.74$, $1.19$, $0.84$, $0.45$~s; pour tout $I$, $T=0.840$~s: l'entrée ne change que l'amplitude, pas le tempo.}
\answer{ex:13:7}{La substitution $s=u-a$ donne la vitesse d'amortissement maximale $a+1/d$ pour $G=e^{-1-ad}/d$; $a\ge33.3$~s$^{-1}$, $G=e^{-2}/d\approx4.5$~s$^{-1}$; plus la co-contraction est forte (plus $a$ est grand), plus il faut prendre $G$ petit.}
\answer{ex:13:8}{Problème ouvert. Par exemple, un seuil dans l'alimentation de la fatigue, $b[r_i-r_0]_+$, donne $I_{\min}=\beta b r_0/(b-\beta+1)\approx0.53$ pour $b=4$, $r_0=0.2$, et une période allant de $1.8$~s pour $I=0.6$ à $0.52$~s pour $I=3$.}
